% =====================================================================
% AI Academy -- National AI Center -- MLE-AI-201 Machine Learning
% HW1 (Weeks 1-2) REPORT TEMPLATE
%
% This template imitates the IEEE two-column conference format WITHOUT
% depending on IEEEtran.cls or any special class file. It uses only the
% article class and packages that ship with every TeX distribution
% (geometry, graphicx, amsmath, amssymb). It compiles with pdflatex or
% tectonic out of the box.
%
% HOW TO USE:
%   1. Fill in the title block (your name + student ID).
%   2. Replace every [FILL IN ...] with your own content.
%   3. Put your plots in ../figures/ and include them with \includegraphics.
%   4. Compile:  pdflatex report_template.tex   (run twice for references)
%      or:       tectonic report_template.tex
%   5. Rename the output to report.pdf before submitting.
% =====================================================================

\documentclass[10pt,twocolumn,letterpaper]{article}

% ---- Minimal, always-available packages only -------------------------
\usepackage[T1]{fontenc}
\usepackage[utf8]{inputenc}
\usepackage[letterpaper,margin=0.75in,columnsep=0.25in]{geometry}
\usepackage{times}          % Times-like font, to mimic the IEEE look
\usepackage{amsmath,amssymb}
\usepackage{graphicx}
\graphicspath{{figures/}{../figures/}}

% ---- Imitate the IEEE section styling by hand (no titlesec) ----------
\renewcommand{\thesection}{\Roman{section}}
\renewcommand{\thesubsection}{\Alph{subsection}}
\makeatletter
% Centered, small-caps section headings: "I. Introduction"
\renewcommand\section{\@startsection{section}{1}{\z@}%
  {3.5ex plus 1ex minus .2ex}{1.8ex plus .2ex}%
  {\centering\normalfont\normalsize\scshape}}
% Italic left-aligned subsection headings: "A. Subsection"
\renewcommand\subsection{\@startsection{subsection}{2}{\z@}%
  {3.0ex plus 1ex minus .2ex}{1.0ex plus .2ex}%
  {\normalfont\normalsize\itshape}}
\makeatother

% Tighter caption look, IEEE-ish
\renewcommand{\figurename}{Fig.}
\setlength{\parindent}{1em}

% =====================================================================
\begin{document}

% ---- TITLE BLOCK (spans both columns) -------------------------------
\twocolumn[
  \begin{center}
    {\LARGE\bfseries KNN, Metrics, and Regression from Scratch\par}
    \vspace{1.2em}
    {\large [FILL IN: Your Name]\par}
    \vspace{0.3em}
    {\normalsize Student ID: [FILL IN] \quad\textbullet\quad Date: [FILL IN]\par}
    \vspace{0.3em}
    {\itshape MLE-AI-201 Machine Learning, Cohort II\\ AI Academy, National AI Center\par}
    \vspace{1.4em}
  \end{center}
]

% ---- ABSTRACT (IEEE style: bold-italic "Abstract" em-dash) ----------
\noindent\textbf{\textit{Abstract}}---\textit{%
[FILL IN: 3--5 sentences. State what you implemented (KNN, four metrics,
linear and logistic regression), the datasets used, your headline
numbers (e.g.\ best $k$ and test F$_1$; linear-regression test $R^2$;
logistic-regression test accuracy), and whether your from-scratch code
matched scikit-learn.]%
}

\vspace{0.6em}
\noindent\textbf{\textit{Index Terms}}---\textit{k-nearest neighbors,
classification metrics, linear regression, logistic regression,
gradient descent, regularization.}

\vspace{0.8em}

% =====================================================================
\section{Introduction}
[FILL IN: One short paragraph. What is the goal of this assignment and
what will the reader find in this report? Mention that every algorithm
is implemented in NumPy and validated against scikit-learn.]

% =====================================================================
\section{KNN and Metrics}
\subsection{KNN implementation}
[FILL IN: Describe how you vectorized the Euclidean distance (no Python
loop over training points). One or two sentences plus the key equation
if helpful.]

The pairwise squared distance can be expanded as
\begin{equation}
  \lVert \mathbf{x}-\mathbf{z}\rVert_2^2
  = \lVert \mathbf{x}\rVert_2^2 - 2\,\mathbf{x}^\top\mathbf{z}
  + \lVert \mathbf{z}\rVert_2^2 .
\end{equation}

\subsection{Metrics implementation}
[FILL IN: Briefly state how you computed accuracy, precision, recall,
and F$_1$ from the confusion-matrix counts, and how you handled the
zero-denominator edge case.]

\subsection{Results}
[FILL IN: Report the metric-vs-$k$ sweep and the best $k$ by F$_1$.
Reference the table and figure below. Confirm the match with
scikit-learn.]

% ---- Example table (plain tabular, no booktabs dependency) ----------
\begin{table}[h]
  \centering
  \caption{Test-set metrics for the best $k$ (breast cancer).}
  \label{tab:knn}
  \begin{tabular}{lcccc}
    \hline
    Model & Acc & Prec & Rec & F$_1$ \\
    \hline
    Your KNN ($k{=}$[..]) & [..] & [..] & [..] & [..] \\
    sklearn KNN           & [..] & [..] & [..] & [..] \\
    \hline
  \end{tabular}
\end{table}

% ---- Example figure -------------------------------------------------
\begin{figure}[h]
  \centering
  % \includegraphics[width=\linewidth]{knn_sweep.pdf}
  \fbox{\parbox[c][3.2cm][c]{0.9\linewidth}{\centering\itshape
    [FILL IN: replace this box with\\ \texttt{\textbackslash includegraphics\{knn\_sweep.pdf\}}]}}
  \caption{Validation F$_1$ as a function of $k$.}
  \label{fig:knn}
\end{figure}

% =====================================================================
\section{Linear and Logistic Regression}
\subsection{Linear regression}
[FILL IN: Describe your batch gradient-descent update, learning rate,
and number of epochs. Report test-set MSE and $R^2$ on the diabetes
dataset and compare with scikit-learn.]

The mean-squared-error objective and its gradient are
\begin{equation}
  J(\mathbf{w},b) = \frac{1}{N}\lVert \mathbf{X}\mathbf{w}+b-\mathbf{y}\rVert_2^2,
  \quad
  \nabla_{\mathbf{w}}J = \frac{2}{N}\mathbf{X}^\top(\mathbf{X}\mathbf{w}+b-\mathbf{y}).
\end{equation}

\subsection{Logistic regression}
[FILL IN: Describe your cross-entropy gradient descent on the breast
cancer dataset. Report test-set accuracy, precision, recall, and F$_1$
(reuse your Part 1 metrics) and compare with scikit-learn.]

The sigmoid and cross-entropy gradient are
\begin{equation}
  \sigma(z)=\frac{1}{1+e^{-z}},
  \quad
  \nabla_{\mathbf{w}}J=\frac{1}{N}\mathbf{X}^\top(\hat{\mathbf{p}}-\mathbf{y}).
\end{equation}

\subsection{Training curves}
[FILL IN: Show and briefly discuss the loss-vs-epoch curves for both
models. A decreasing curve confirms the learning rate is sensible.]

% =====================================================================
\section{Bonus: Regularization (optional)}
[FILL IN, only if attempted. State clearly which bonus(es) you did.
For L1: show that a large $\lambda$ drives some weights toward zero.
For L2: show that weights shrink smoothly. Report how the weight norm
and test performance change with $\lambda$.]

% =====================================================================
\section{Conclusion}
[FILL IN: 2--3 sentences. What worked, what surprised you, and where
your from-scratch results differed from scikit-learn (and why).]

% =====================================================================
\section*{AI-Tool Usage Declaration}
[FILL IN: Name any AI tool used, what you used it for, which parts it
affected, and how you checked the output. Write "No AI tools used." if
that is the case. This declaration is required by the syllabus.]

% =====================================================================
\begin{thebibliography}{9}
\bibitem{islp} G.~James, D.~Witten, T.~Hastie, R.~Tibshirani, and J.~Taylor,
  \textit{An Introduction to Statistical Learning with Applications in Python}. Springer, 2023.
\bibitem{esl} T.~Hastie, R.~Tibshirani, and J.~Friedman,
  \textit{The Elements of Statistical Learning}, 2nd ed. Springer, 2009.
\bibitem{sklearn} F.~Pedregosa et al., ``Scikit-learn: Machine learning in Python,''
  \textit{JMLR}, vol.~12, pp.~2825--2830, 2011.
\end{thebibliography}

\end{document}
